\begin{table}[H]
\begin{center}
\begin{threeparttable}
\captionsetup{margin=0.5cm}
\caption{Regional acceptance of LGB and transgender children in Chile.}
\label{tab:enssex_acceptance_region}
\footnotesize
\centering
\begin{tabular}{lcccc}
\toprule
Region & \multicolumn{2}{c}{LGB child} & \multicolumn{2}{c}{Trans child} \\
\cmidrule(lr){2-3}\cmidrule(lr){4-5}
 & Observations & Acceptance (\%) & Observations & Acceptance (\%) \\
\midrule
Tarapacá & 607 & 75.96 & 607 & 71.19 \\
Antofagasta & 883 & 82.44 & 881 & 81.43 \\
Atacama & 524 & 81.33 & 525 & 78.34 \\
Coquimbo & 1169 & 84.41 & 1169 & 83.16 \\
Valparaíso & 1680 & 87.38 & 1679 & 85.59 \\
O'Higgins & 1393 & 81.61 & 1396 & 76.43 \\
Maule & 1431 & 80.34 & 1429 & 76.76 \\
Biobío & 1834 & 73.82 & 1839 & 70.38 \\
La Araucanía & 1475 & 78.17 & 1477 & 71.46 \\
Los Lagos & 1491 & 78.14 & 1476 & 72.40 \\
Aysén & 495 & 77.82 & 492 & 75.20 \\
Magallanes & 437 & 76.74 & 438 & 70.87 \\
Metropolitana & 4361 & 82.98 & 4352 & 78.67 \\
Los Ríos & 638 & 82.80 & 637 & 79.70 \\
Arica y Parinacota & 518 & 77.84 & 519 & 71.24 \\
Ñuble & 1145 & 84.48 & 1144 & 82.79 \\
\midrule
National & 20081 & 81.84 & 20060 & 78.01 \\
\bottomrule
\end{tabular}
\begin{tablenotes}
\item Notes: Acceptance rates are weighted using the ENSSEX 2022--2023 calibrated person weight.\item The number of observations corresponds to the unweighted number of valid responses in each region.
\item Source: Own elaboration based on ENSSEX 2022--2023.
\end{tablenotes}
\end{threeparttable}
\end{center}
\end{table}
